\newpage
\phantomsection
\addcontentsline{toc}{section}{致谢}
\section*{致~~谢}
\vspace*{-2mm}
% 2026: 致谢正文 小四 固定25pt 与正文一致
\zihao{-4}\setlength{\baselineskip}{25pt}

时光匆匆，四年的本科生活即将画上句号。回首这段旅程，有太多的感激想要表达。

首先，我要衷心感谢我的指导教师王红老师。从选题方向的确定到论文框架的搭建，从实验方案的讨论到文稿的反复修改，王老师始终给予我耐心细致的指导。每当我在研究中遇到瓶颈、感到迷茫时，王老师总能用清晰的思路帮我拨开迷雾，让我重新找到前进的方向。王老师严谨的治学态度和对细节的执着追求，是我在整个毕业设计过程中最宝贵的收获。

感谢学院的各位老师，四年来你们在课堂上传授的知识，为我打下了坚实的专业基础。正是信号处理、机器学习、软件工程等课程的积累，让我有能力完成这样一个跨领域的课题。

感谢我的同学和室友们。深夜实验室里的相互鼓励、食堂里关于技术方案的激烈讨论、答辩前的模拟演练——这些看似平常的日子，现在回想起来都格外珍贵。感谢你们在我焦虑时给我力量，在我懈怠时督促我前行。

感谢开源社区的无私奉献。本文的研究建立在 MNE-Python、scikit-learn、PyTorch 等优秀开源项目之上，BCI Competition IV 数据集的公开共享也为本研究提供了不可或缺的实验基础。正是这种开放共享的精神，让一个本科生也有机会站在巨人的肩膀上做一点有意义的探索。

最后，我要感谢我的父母。感谢你们二十多年来无条件的爱与支持，感谢你们尊重我每一次选择、包容我每一次犯错。你们是我最坚实的后盾，也是我不断向前的动力。

本科四年，是结束，也是开始。愿此去繁花似锦，再相逢依然如故。






\label{thanks}  % 请勿删除此行



